\begin{proofbox}
\par\noindent\begin{minipage}{\linewidth}
\textbf{\textcolor{CProof}{证明目标与要注意的地方}}
已知 $a>b>0$，证明椭圆
\[
\frac{x^2}{a^2}+\frac{y^2}{b^2}=1
\]
上所有点都能写作 $P(a\cos\theta,b\sin\theta)$，反过来这样写出的点也一定在椭圆上。\textbf{只证明“代入后等于 1”还不够}，因为那只能说明写出的点都对，不能说明椭圆上的点一个没漏。

\medskip
\end{minipage}\par
\par\noindent\begin{minipage}{\linewidth}
\textbf{\textcolor{CProof}{方法一：门槛最低——初中直角三角形、圆和压扁图形}}

\begin{center}
\begin{tikzpicture}[x=1cm,y=1cm,font=\small,>=Stealth]
\draw[rounded corners=3pt,gray!45] (-6.3,-4.05) rectangle (6.3,4.05);
\node[font=\small\bfseries] at (0,3.5999999999999996) {方法一：压缩与还原互为逆操作};
\draw[->,gray] (-3.5,0)--(3.7,0) node[right] {$x$};\draw[->,gray] (0,-3.15)--(0,3.15) node[right] {$y$};\node[below left] at (0,0) {$O$};\draw[orange!80!black,dashed,thick] (0,0) circle (3);\draw[CStatement,very thick] (0,0) ellipse (3 and 1.8);
\draw[orange!80!black,thick] (0,0)--(1.8,2.4)--(1.8,0);
\draw (1.65,0)--(1.65,.15)--(1.8,.15);
\draw[->,CProof,thick] (1.95,2.36)--(1.95,1.48);
\draw[->,teal!80!black,thick] (1.65,1.48)--(1.65,2.36);
\fill[orange!80!black] (1.8,2.4) circle (2pt) node[above right] {$Q$};
\fill[CStatement] (1.8,1.44) circle (2pt);\draw[gray] (1.8,1.44)--(3,.9);\node at (3.3,.9) {$P$};
\node[below] at (1.8,0) {$a\cos\theta$};
\draw[->,orange!80!black] (.6,0) arc (0:53.130102:.6);
\node at (.85,.4) {$\theta$};
\node[align=left,anchor=west] at (3.5,1.8) {$Q\to P$\\纵坐标乘 $b/a$};
\node[align=left,anchor=west] at (3.5,-1.5) {$P\to Q$\\纵坐标乘 $a/b$};
\node[align=right,anchor=east] at (-3.0,-1.3) {直角三角形给出\\圆上的坐标};

\end{tikzpicture}
\end{center}

\end{minipage}\par
\begin{description}[style=nextline,leftmargin=2.8em,labelsep=0.5em,itemsep=0.55em]
\item[\textbf{步骤 1：先画一个半径为 $a$ 的圆。}]
圆的方程是 $x^2+y^2=a^2$。取圆上一点 $Q$，从 $Q$ 向 $x$ 轴作垂线。如果在第一象限，按直角三角形里正弦、余弦的定义，
\[
x_Q=a\cos\theta,\qquad y_Q=a\sin\theta.
\]
其他象限把横、纵坐标的正负号对应处理，仍然是这个式子。
\item[\textbf{步骤 2：把 $Q$ 上下压到 $P$。}]
横坐标保持不变，纵坐标变成原来的 $\frac ba$，所以
\[
x_P=x_Q=a\cos\theta,\qquad
y_P=\frac ba y_Q=b\sin\theta.
\]
\textbf{这一步的意义：} 横向最大长度仍是 $a$，纵向最大长度从 $a$ 变成 $b$。
\item[\textbf{步骤 3：确认 $P$ 在椭圆上。}]
\[
\frac{x_P^2}{a^2}+\frac{y_P^2}{b^2}
=\cos^2\theta+\sin^2\theta=1.
\]
\item[\textbf{步骤 4：倒着走，确认没有漏点。}]
任取椭圆上的点 $P(x,y)$，把它的纵坐标放大 $\frac ab$ 倍，得 $Q(x,\frac ab y)$。因为
\[
x^2+\left(\frac ab y\right)^2
=a^2\left(\frac{x^2}{a^2}+\frac{y^2}{b^2}\right)=a^2,
\]
所以 $Q$ 在原来的圆上，必能找到一个 $\theta$，使 $Q=(a\cos\theta,a\sin\theta)$。再上下压扁便回到 $P=(a\cos\theta,b\sin\theta)$。
\end{description}
\textbf{结论：} 圆上的全部点上下压扁，恰好得到椭圆上的全部点。

\medskip
\par\noindent\begin{minipage}{\linewidth}
\textbf{\textcolor{CProof}{方法二：考场速度最快——两个分数拼成单位圆}}

\begin{center}
\begin{tikzpicture}[x=1cm,y=1cm,font=\small,>=Stealth]
\draw[rounded corners=3pt,gray!45] (-6.3,-3.45) rectangle (6.3,3.45);
\node[font=\small\bfseries] at (0,3.0) {方法二：两个分数就是单位圆的两个坐标};

\draw[->,gray] (-2.3,0)--(2.5,0) node[right] {$u=x/a$};
\draw[->,gray] (0,-2.3)--(0,2.5);\node[anchor=west] at (.25,2.4) {$v=y/b$};
\draw[orange!80!black,thick] (0,0) circle (2);
\draw[CStatement,thick] (0,0)--(1.2,1.6);
\draw[dashed] (1.2,0)--(1.2,1.6)--(0,1.6);
\draw (1.05,0)--(1.05,.15)--(1.2,.15);
\draw[->] (.65,0) arc (0:53.130102:.65);
\node at (.9,.45) {$\theta$};\node[below left] at (0,0) {$O$};
\node[below right] at (2.06,-.06) {$1$};\node[above left] at (0,2) {$1$};
\node[below] at (1.2,0) {$\cos\theta$};\node[left] at (0,1.6) {$\sin\theta$};
\fill[CStatement] (1.2,1.6) circle (2pt) node[above right] {$U(x/a,y/b)$};
\node[anchor=east,align=right] at (-2.8,-1.3) {归一化后的平面\\$u^2+v^2=1$};
\node[anchor=west,align=left] at (3,-1.3) {半径为 $1$\\$u=\cos\theta$\\$v=\sin\theta$};

\end{tikzpicture}
\end{center}

\end{minipage}\par
\begin{description}[style=nextline,leftmargin=2.8em,labelsep=0.5em,itemsep=0.45em]
\item[\textbf{步骤 1：先看椭圆方程。}]
它说的是 $\left(\frac xa\right)^2+\left(\frac yb\right)^2=1$。
\item[\textbf{步骤 2：想到半径为 $1$ 的圆。}]
在单位圆上，点的两个坐标可写成 $\cos\theta,\sin\theta$。因此能取一个角，使
\[
\frac xa=\cos\theta,\qquad \frac yb=\sin\theta.
\]
这一步有依据：$(\frac xa,\frac yb)$ 的两个坐标平方和等于 $1$，确实位于单位圆上。
\item[\textbf{步骤 3：乘回去。}]
得到 $x=a\cos\theta,y=b\sin\theta$。反过来代入椭圆方程，又有 $\cos^2\theta+\sin^2\theta=1$。
\end{description}
\textbf{考场识别信号：} 方程中出现\textbf{两个平方相加等于 $1$}，就试着把它们看成单位圆上一个点的横、纵坐标。

\medskip
\par\noindent\begin{minipage}{\linewidth}
\textbf{\textcolor{CProof}{方法三：步骤最少——只选一个坐标，再把另一个补出来}}

\begin{center}
\begin{tikzpicture}[x=1cm,y=1cm,font=\small,>=Stealth]
\draw[rounded corners=3pt,gray!45] (-6.3,-3.6) rectangle (6.3,3.6);
\node[font=\small\bfseries] at (0,3.15) {方法三：横坐标相同，纵坐标的正负决定选哪个角};

\draw[->,gray] (-3.5,0)--(3.7,0) node[right] {$x$};
\draw[->,gray] (0,-2.3)--(0,2.5) node[right] {$y$};
\draw[CStatement,very thick] (0,0) ellipse (3 and 1.8);
\draw[dashed] (1.8,-1.44)--(1.8,1.44);
\draw[teal!80!black,thick] (0,0)--(1.8,1.44);
\draw[CTrap,thick] (0,0)--(1.8,-1.44);
\fill[teal!80!black] (1.8,1.44) circle (2pt) node[above right] {$P_+(x,b\sin\theta_0)$};
\fill[CTrap] (1.8,-1.44) circle (2pt) node[below right] {$P_-(x,-b\sin\theta_0)$};
\node[anchor=west] at (3.5,.75) {$x=a\cos\theta_0$};
\node[anchor=east,align=right] at (-3.0,1.5) {上半部：$\theta_0$\\$0<\theta_0<\pi$};
\node[anchor=east,align=right] at (-3.0,-1.5) {下半部：$2\pi-\theta_0$\\余弦相同\\正弦相反};
\node[below left] at (0,0) {$O$};
\node at (0,-2.8) {$x=\pm a$ 时上下两点合一，$y=0$};

\end{tikzpicture}
\end{center}

\end{minipage}\par
\begin{description}[style=nextline,leftmargin=2.8em,labelsep=0.5em,itemsep=0.45em]
\item[\textbf{步骤 1：确定横坐标可以怎样写。}]
从 $\frac{x^2}{a^2}+\frac{y^2}{b^2}=1$ 得 $0\le\frac{x^2}{a^2}\le1$，所以 $-a\le x\le a$。于是总可以找一个角 $\theta$，使 $x=a\cos\theta$。
\item[\textbf{步骤 2：用方程补出纵坐标的大小。}]
\[
\frac{y^2}{b^2}=1-\cos^2\theta=\sin^2\theta,
\qquad |y|=b|\sin\theta|.
\]
\item[\textbf{步骤 3：别忘了纵坐标有正负。}]
同一个 $\cos\theta$（除左右端点外）通常对应圆上、下两个角。若 $y>0$，选 $\sin\theta>0$ 的角；若 $y<0$，选 $\sin\theta<0$ 的角；若 $y=0$，选 $\theta=0$ 或 $\pi$。因此总可以让 $y=b\sin\theta$。
\item[\textbf{步骤 4：代回检查。}]
直接得到 $\frac{(a\cos\theta)^2}{a^2}+\frac{(b\sin\theta)^2}{b^2}=1$。
\end{description}
\textbf{这条路的好处：} 从横坐标范围直接出发；\textbf{最容易漏掉的地方：} 从 $y^2=b^2\sin^2\theta$ 直接写成 $y=b\sin\theta$，却不交代怎样选角使符号一致。

\medskip
\par\noindent\begin{minipage}{\linewidth}
\textbf{\textcolor{CProof}{方法四：常规证明——先验证，再证明每个点都可取到}}

\begin{center}
\begin{tikzpicture}[x=1cm,y=1cm,font=\small,>=Stealth]
\draw[rounded corners=3pt,gray!45] (-6.3,-3.6) rectangle (6.3,3.6);
\node[font=\small\bfseries] at (0,3.15) {方法四：从任意 P 出发，反找单位圆上的 U};

\begin{scope}[shift={(-2.9,0)},scale=.72]
\draw[->,gray] (-3.4,0)--(3.5,0) node[right] {$x$};
\draw[->,gray] (0,-2.1)--(0,2.2) node[right] {$y$};
\draw[CStatement,thick] (0,0) ellipse (3 and 1.8);
\fill[CStatement] (-1.8,-1.44) circle (2.5pt) ;\node at (-1.8,-2.3) {$P(x_0,y_0)$};
\node at (0,2.85) {任取椭圆上一点};
\end{scope}
\begin{scope}[shift={(3.2,0)}]
\draw[->,gray] (-1.7,0)--(1.8,0) node[right] {$u$};
\draw[->,gray] (0,-1.7)--(0,1.8) node[right] {$v$};
\draw[orange!80!black,thick] (0,0) circle (1.4);
\draw[CStatement,thick] (0,0)--(-.84,-1.12);
\draw[->] (.5,0) arc (0:233.130102:.5);
\node at (-.75,.55) {$\theta$};
\fill[CStatement] (-.84,-1.12) circle (2pt) ;\node at (-.5,-1.95) {$U(x_0/a,y_0/b)$};
\node at (0,2.05) {必在单位圆上};
\end{scope}
\draw[->,CProof,thick] (-.4,.6)--(1.1,.6);
\node[above,align=center] at (.35,.7) {分别除以\\$a,b$};
\draw[->,teal!80!black,thick] (1.1,-.5)--(-.4,-.5);
\node[below] at (.35,-.6) {乘回 $a,b$};
\node at (0,-2.65) {图取第三象限：其他象限及轴上点同样能找到 $\theta\in[0,2\pi)$};

\end{tikzpicture}
\end{center}

\end{minipage}\par
\begin{description}[style=nextline,leftmargin=2.8em,labelsep=0.5em,itemsep=0.55em]
\item[\textbf{步骤 1：验证设出来的点。}]
任取 $0\le\theta<2\pi$，代入可得
\[
\begin{aligned}
\frac{(a\cos\theta)^2}{a^2}+\frac{(b\sin\theta)^2}{b^2}
&=\cos^2\theta+\sin^2\theta\\
&=1.
\end{aligned}
\]
因此设出的每个点都在椭圆上。
\item[\textbf{步骤 2：反向任取椭圆上一点 $P(x_0,y_0)$。}]
由椭圆方程可知
\[
\left(\frac{x_0}{a}\right)^2+\left(\frac{y_0}{b}\right)^2=1.
\]
把数对 $\left(\frac{x_0}{a},\frac{y_0}{b}\right)$ 画到平面中，它距原点的距离为 $1$，所以就在单位圆上。
\item[\textbf{步骤 3：找到对应的角。}]
从 $x$ 轴正方向转到这个单位圆上的点，总有一个角 $\theta\in[0,2\pi)$，使
\[
\cos\theta=\frac{x_0}{a},\qquad \sin\theta=\frac{y_0}{b}.
\]
于是 $P=(a\cos\theta,b\sin\theta)$，也就是说椭圆上的任何一点都能取到。
\item[\textbf{步骤 4：两边合在一起。}]
已证明“设出来的一定在椭圆上”和“椭圆上的都能设出来”，所以两种表示完全对应。
\end{description}

\medskip
\textbf{\textcolor{CProof}{四种方法怎样选？}}
方法一适合\textbf{第一次理解为什么这样设}；方法二适合\textbf{考场快速识别}；方法三适合\textbf{从代数方程自己推出来}；方法四适合\textbf{写完整的书面证明}。它们并不是四个彼此无关的结论，而是同一件事从不同起点来说明。
\end{proofbox}
